\newtoggle{withbonus} \toggletrue{withbonus} % set true to show bonus points in grade table

% setting underlying course name (Grundlagenfach Mathematik, §-Unterricht, \ldots)
\newcommand{\mycourse}{Grundlagenfach Informatik}
% name of the class
\newcommand{\myclass}{Klasse 2d}
% set date
\newcommand{\mydate}{Donnerstag, 09.04.2026}

\title{\textbf{Hardware, Endliche Automaten, Laufzeitkomplexität}}
\author{Thomas Graf\\
	\mycourse, \myclass}
\date{\mydate}
\maketitle
\thispagestyle{headandfoot}

\begin{center}
	\fbox{\parbox{140mm}{\centering
			\begin{itemize}
				\item Dauer der Prüfung: $45$ Minuten
				      %\item \textbf{Zeige in jeder Aufgabe unbedingt Deine Schritte!}
				\item Erlaubte Hilfsmittel:
				      \begin{itemize}
					      \item Schreibzeug
					            %\item Taschenrechner
				      \end{itemize}
			\end{itemize}}}
\end{center}
\vspace{10mm}

\begin{center}
	Vorname: \rule{45mm}{0.8pt}\qquad Nachname: \rule{45mm}{0.8pt}
\end{center}
\vspace{20mm}

\begin{center}
	\iftoggle{withbonus} {
		\combinedgradetable[h][questions]
	} {
		\gradetable[h][questions]
	}
\end{center}
\vspace{15mm}
\begin{center}
	Note:
	\vspace{10mm}

	\rule{8mm}{1.2pt}
\end{center}
%\thispagestyle{empty}
\clearpage

\begin{questions}
    \question
    \begin{parts}
        \part[1] Wie lässt sich die Logikfunktion für das Summenbit ($\Sigma$) eines Volladdierers am effizientesten mit \textsc{XOR}-Gattern beschreiben? Dabei sind $A$ und $B$ die beiden zu addierenden Bits und $C_{in}$ das Carry-In-Bit und $\oplus$ bezeichnet die \textsc{XOR}-Operation.
            \begin{checkboxes}
                \correctchoice $\Sigma = A \oplus B \oplus C_{in}$
                \choice $\Sigma = (A\;\textsc{AND}\; B) \oplus C_{in}$
                \choice $\Sigma = A\;\textsc{AND}\; B\; \textsc{OR}\; C_{in}$
                \choice keine der obigen Antworten ist richtig
            \end{checkboxes}
        \part[1] Welche logische Funktion wird durch zwei in Reihe (hintereinander) geschaltete Schalter $A$ und $B$ realisiert, damit eine LED leuchtet?
            \begin{checkboxes}
                \choice \textsc{XOR}
                \choice \textsc{OR}
                \correctchoice \textsc{AND}
            \end{checkboxes}
        \end{parts}
        \droptotalpoints


    \question[2]{\textbf{Anzahl der Elemente in einer Sprache bestimmen}}

	Es sei $\Sigma = \lrc{a, b, c, d}$ ein Alphabet und $n$ eine gegebene natürliche Zahl. Bestimmen Sie die Anzahl der verschiedenen Wörter der Länge $n$ über $\Sigma$, welche mindestens ein $a$ enthalten.
	\begin{solution}
		Insgesamt gibt es genau $\abs{\Sigma}^n = 4^n$ Wörter der Länge $n$ über $\Sigma$. Davon enthalten genau $3^n$ Wörter kein $a$ (da für jede Stelle nur die 3 Buchstaben $\{b, c, d\}$ zur Auswahl stehen). Die gesuchte Anzahl der Wörter mit mindestens einem $a$ ist somit $4^n - 3^n$.
	\end{solution}
\vspace{5cm}

\clearpage
	\question{\textbf{Endlichen Automaten zeichnen}}

	Gegeben ist die Sprache
	\begin{align*}
		L_1 := \setcm{0x11y}{x,y\in\lrc{0,1}^*}.
	\end{align*}
	\begin{parts}
		\part[1] Nennen Sie ein kürzestes Wort in $L_1$.
		\begin{solution}
			Das Wort $011$ ist \textbf{das} (eindeutige) kürzeste Wort in $L_1$.
		\end{solution}
		\vspace{5cm}
        \part[2] Entwerfen Sie einen endlichen Automaten (in Diagrammform), welcher $L_1$ akzeptiert.
		\begin{solution}
			\begin{figure}[H]
				\centering
				\begin{tikzpicture}[->,>=stealth',shorten >=1pt,node distance=3.2cm,semithick]
					\tikzstyle{every state}=[fill=green!10,draw=black,text=black,minimum size=1.4cm]

					\node[initial,state] (q0) {$q_{0}$};
					\node[state] (q1) [above of=q0] {$q_{1}$};
					\node[state] (q2) [right of=q0] {$q_{2}$};
					\node[state] (q3) [right of=q2] {$q_{3}$};
					\node[state,accepting] (q4) [right of=q3] {$q_{4}$};

					\path (q0) edge [left] node {$1$} (q1)
					(q0) edge [below] node {$0$} (q2)

					(q1) edge [loop above] node {$0,1$} (q1)

					(q2) edge [below] node {$1$} (q3)
					(q2) edge [loop above] node {$0$} (q2)

					(q3) edge [above, bend right] node {$0$} (q2)
					(q3) edge [below] node {$1$} (q4)

					(q4) edge [loop above] node {$0,1$} (q4);
				\end{tikzpicture}
				\caption{Endlicher Automat (in Diagrammform) für die Sprache $L_1$}
				\label{fig:L1}
			\end{figure}
		\end{solution}
	\end{parts}
	\droptotalpoints

    \clearpage
	\question{\textbf{\red{Nicht reguläre} Sprache}}

	Gegeben sei die Sprache
	\begin{align*}
		L_2 := \setcm{w\$\$\$w^2}{w\in\lrc{0, 1}^*}
	\end{align*}
	über dem Alphabet $\lrc{0, 1, \$}$.

	(Zur Vermeidung von Missverständnissen: Für beispielsweise $w = 011$ ist $w^2 = (011)^2 = 011011$.)

	\begin{parts}
		\part[1] Nenne genau drei verschiedene Wörter, die in $L_2$ liegen.
		\begin{solution}
			\begin{align*}
				\$\$\$, 0\$\$\$00, 10\$\$\$1010
			\end{align*}
		\end{solution}
		\vspace{3cm}
		\part[3] Zeige, dass $L_2$ \red{nicht regulär} ist.

		\begin{solution}
			Angenommen, $L_2$ sei regulär. Dann existiert ein endlicher Automat
			\begin{align*}
				A = (Q, \{0, 1, \$\}, \delta, q_0, F)
			\end{align*}
			und $A$ akzeptiert $L_2$. Sei jetzt $m := \abs{Q}$, also die Anzahl der verschiedenen Zustände des endlichen Automaten
			$A$. Wir betrachten die $m+1$ verschiedenen Wörter
			\begin{align*}
				0^1, 0^2, \ldots, 0^{m+1}.
			\end{align*}
			Dies sind mehr Worte (genau ein Wort mehr), als $A$ Zustände hat. Nach dem Schubfachprinzip (pigeonhole principle) existieren
			$i,j\in\lrc{1,2,\ldots,m+1}$ mit $i<j$, sodass
			\begin{align*}
				\hat{\delta}\lr{q_0,0^i} = \hat{\delta}\lr{q_0,0^j}.
			\end{align*}
			Nach der im Unterricht bewiesenen Aussage gilt
			\begin{align*}
				0^iz\in L_2 \iff 0^jz\in L_2
			\end{align*}
			für alle $z\in\{0, 1, \$\}^*$. Doch dies gilt nicht, da zum Beispiel für die Wahl $z:=\$\$\$0^i0^i$
			\begin{align*}
				0^iz = 0^i\$\$\$0^i0^i\in L_2 \quad\text{aber}\quad 0^jz = 0^j\$\$\$0^i0^i\notin L_2
			\end{align*}
			gilt, da $0^i0^i$ zwar das Quadrat von $0^i$ ist, aber nicht das Quadrat von $0^j$ (da $i < j$). Damit war unsere ursprüngliche Annahme falsch und es existiert kein
			endlicher Automat für die Sprache $L_2$.
		\end{solution}
	\end{parts}
	\droptotalpoints

    \clearpage
        \question Bestimme jeweils die Worst-Case-Laufzeit der folgenden Algorithmen in Abhängigkeit von $n\in\N$.
        \begin{parts}
            \part[1] \leavevmode
            \begin{lstlisting}[language=Python]
                result = 0
                for i in range(n):
                    for j in range(i, n):
                        result += i + j
                \end{lstlisting}
            \begin{checkboxes}
                \choice linear
                \choice proportional zu $\sqrt{n}$
                \choice logarithmisch
                \correctchoice quadratisch
                \choice proportional zu $n \log(n)$
                \choice keine der obigen Antworten ist richtig
            \end{checkboxes}

            \part[1] \leavevmode
            \begin{lstlisting}[language=Python]
                result = 0
                i = 1
                while i <= n:
                    result += i
                    i *= 2
            \end{lstlisting}
            \begin{checkboxes}
                \correctchoice logarithmisch
                \choice proportional zu $\sqrt{n}$
                \choice linear
                \choice quadratisch
                \choice proportional zu $n \log(n)$
                \choice keine der obigen Antworten ist richtig
            \end{checkboxes}
            \part[1] \leavevmode
            \begin{lstlisting}[language=Python]
                result = 0
                for i in range(n):
                    j = 0
                    while j < n:
                        result += 1
                        j += 1
            \end{lstlisting}
            \begin{checkboxes}
                \choice linear
                \choice proportional zu $\sqrt{n}$
                \choice logarithmisch
                \correctchoice quadratisch
                \choice proportional zu $n \log(n)$
                \choice keine der obigen Antworten ist richtig
            \end{checkboxes}

            \part[1] \leavevmode
            \begin{lstlisting}[language=Python]
                result = 0
                i = n
                while i >= 1:
                    for j in range(i):
                        result += 1
                    i //= 2
            \end{lstlisting}
            \begin{checkboxes}
                \choice proportional zu $n \log(n)$
                \choice proportional zu $\sqrt{n}$
                \choice logarithmisch
                \choice quadratisch
                \correctchoice linear
                \choice keine der obigen Antworten ist richtig
            \end{checkboxes}
        \end{parts} \droptotalpoints

\end{questions}